% ===========================================================================
% Capitolo 4 — Implementazione
% Tesi triennale: Francesco Zamar, UniTS, A.A. 2025-2026
% Relatore: Prof. Andrea De Lorenzo
% Single source of truth per numeri e macro: macros_finale.tex
% ===========================================================================

\chapter{Implementazione}
\label{ch:implementazione}

L'architettura del \Cref{ch:progettazione} è una scelta di responsabilità:
cosa decide il layer simbolico, cosa spiega il layer neurale, come si
contrattano via REST. Questo capitolo mostra come quelle responsabilità si
traducono in codice eseguibile, riproducibile e testato, e quali compromessi
tecnologici si sono dovuti accettare per rispettare i vincoli di
determinismo, riproducibilità e deployment locale fissati a progetto. Le
scelte concrete (Python 3.12, FastAPI, ChromaDB embedded, Ollama, Pydantic)
non sono frutto di una rassegna comparativa esaustiva: sono le opzioni
che la combinazione di richiesta clinica e ambiente di calcolo disponibile
restringeva a un'unica soluzione naturale. La trattazione segue l'ordine
logico dei componenti, dallo stack tecnologico alla demo interattiva,
ricalcando la mappa concettuale del capitolo precedente per mantenerne la
leggibilità.

% ============================================================
\section{Stack tecnologico e organizzazione del codice}
\label{sec:stack}

L'intero sistema è scritto in Python~3.12, scelto per la maturità
dell'ecosistema scientifico, la tipizzazione statica progressiva introdotta
dalle versioni recenti e la facilità di integrazione con le librerie di
machine learning necessarie alla pipeline RAG (\emph{Retrieval-Augmented Generation}). Il codice di produzione ammonta
a circa 10\,000 righe, distribuite fra quattro moduli principali contenuti
nella directory \texttt{Note\_AIFA/} del repository.

Il modulo \texttt{aifa\_rule\_engine/} implementa tutta la logica simbolica
del motore: circa 3\,500 righe tra regole, valutatore, logica trivalente
e API FastAPI~\cite{fastapi} esposta sulla porta 8000. Il modulo
\texttt{rag\_pipeline/} comprende l'ingestione, il retriever ChromaDB, il
costruttore di prompt, i backend LLM (\emph{Large Language Model}) e l'orchestratore FastAPI sulla porta
8001, per circa 2\,500 righe. Il modulo \texttt{evaluation/} raccoglie i
runner di metriche, i report JSON e gli script overnight (circa 3\,000
righe), mentre il modulo \texttt{demo/} realizza l'interfaccia Streamlit in
circa 1\,000 righe.

Il modello linguistico locale è \llmmodel, servito tramite Ollama~\cite{ollama2024}
in un container dedicato; gli embedding semantici sono generati da
\embedmodel~\cite{reimers2019sbert} e il reranking utilizza
\reranker~\cite{nickprock2023}. La persistenza semantica è affidata a
ChromaDB~\cite{chromadb}, mentre la validazione degli schemi JSON è delegata
a Pydantic~\cite{pydantic}.

I tre servizi (rule engine, orchestratore e Ollama) sono orchestrati da
Docker Compose su una rete interna isolata dal traffico esterno. La
comunicazione avviene su socket TCP locali; l'orchestratore può caricare il
rule engine direttamente in-process (variabile d'ambiente
\texttt{AIFA\_RULES\_DIR}) per eliminare la latenza HTTP nei test di
integrazione. La suite di test conta \testcountall{} casi verdi: \testcountengine{}
per il rule engine e \testcountorch{} per l'orchestratore RAG.

% ============================================================
\section{Rule engine: implementazione}
\label{sec:rule-engine-impl}

Il rule engine deterministico è alla versione \engineversion{} ed è scritto
in Python puro, senza dipendenze da framework di reasoning esterni, con
l'obiettivo di mantenere ogni passo computazionale ispezionabile riga per
riga. Il catalogo delle regole è codificato in file YAML per ciascuna nota
(\texttt{rules/nota\_01/}, \texttt{nota\_13/}, \texttt{nota\_66/},
\texttt{nota\_97/}) e caricato al momento dell'avvio del servizio dal
modulo \texttt{engine/rule\_loader.py}. Ogni documento YAML è validato contro
uno schema Pydantic che garantisce la presenza dei campi obbligatori prima
che il motore entri in produzione.

\subsection{Schema YAML delle regole}

Ogni voce del catalogo specifica un identificatore univoco, il tipo logico
della regola, la versione di schema, la descrizione in italiano, l'ordine di
valutazione, un'\emph{anchor} normativa che punta alla pagina e alla sezione
del PDF AIFA (Agenzia Italiana del Farmaco) di riferimento, e la condizione booleana strutturata. Il seguente
estratto è tratto da \texttt{rules/nota\_13/rules.yaml} e illustra una regola
di inclusione e una di eccezione:

\begin{verbatim}
- rule_id: N13_INCL_001
  rule_type: INCLUSION
  nota: '13'
  version: 3.3.0
  description_it: Dieta appropriata seguita per almeno 3 mesi.
  evaluation_order: 10
  normative_anchor:
    pdf_file: nota-13.pdf
    page: 1
    section: Condizioni di prescrizione
    excerpt: "dieta, seguita per almeno tre mesi"
  condition:
    operator: IS_TRUE
    var: dieta_seguita_almeno_3_mesi
\end{verbatim}

\begin{verbatim}
- rule_id: N13_EXCEPT_001
  rule_type: EXCEPTION
  nota: '13'
  version: 3.3.0
  outcome_if_true: BYPASS
  bypasses:
    - N13_INCL_001
  condition:
    operator: AND
    operands:
      - {operator: IS_TRUE, var: rischio_molto_alto}
      - {operator: GTE, var: ldl, value: 70}
\end{verbatim}

Il campo \texttt{normative\_anchor} consente al motore di tracciare ogni
decisione fino alla pagina esatta del documento normativo, soddisfacendo il
requisito di spiegabilità citazionale discusso nel \Cref{ch:progettazione}.

\subsection{Tipi di regole e fasi di valutazione}

Il valutatore implementa dieci fasi sequenziali. La fase SCOPE verifica
l'applicabilità della nota al caso clinico; la fase EXCEPTION identifica
eventuali bypass delle regole di inclusione prima di valutarle; la fase
EXCL\_HARD applica le esclusioni assolute con fail-fast su \texttt{TRUE}; la
fase INCLUSION valuta i prerequisiti del trattamento con fail-fast su
\texttt{FALSE}; la fase PATHWAY verifica il percorso terapeutico; le fasi
GUIDANCE\_DOSE, GUIDANCE\_PREF e GUIDANCE\_WARN producono i flag clinici
accessori senza influire sulla decisione binaria. La fase di finalizzazione
aggrega le evidenze, applica l'invariante di sicurezza e determina l'esito.

\subsection{Implementazione della logica Kleene}

La semantica trivalente è stata richiamata nella \Cref{subsec:kleene};
qui se ne descrive soltanto la codifica concreta. Il modulo
\texttt{logic/three\_valued.py} esporta l'enum \texttt{TruthValue} con i
valori \texttt{TRUE}, \texttt{FALSE} e \texttt{UNKNOWN}, e implementa le
tavole di verità esplicitamente mediante confronti con \texttt{None}: ogni
comparazione che coinvolge un campo non valorizzato restituisce
\texttt{UNKNOWN}. La logica non è delegata a un solutore di terze parti,
garantendo il pieno controllo sul comportamento in presenza di dati
incompleti e la copertura analitica per test di unità.

L'operatore esteso \texttt{SCORE\_RANGE\_GTE} implementa l'\emph{interval
arithmetic} per i punteggi clinici compositi come il CHA\textsubscript{2}DS\textsubscript{2}-VASc: i
componenti ignoti contribuiscono con il minimo pari a~0 e il massimo pari al
loro peso massimo, e il confronto con la soglia $k$ restituisce \texttt{TRUE}
se il punteggio minimo supera già $k$, \texttt{FALSE} se il massimo non la
raggiunge, \texttt{UNKNOWN} negli altri casi~\cite{deStefano2014cha2ds2}.

La decisione \decision{NON\_DETERMINABILE} viene emessa quando almeno un
campo in \texttt{missing\_fields\_coverage} è determinante per l'esito e il
motore non può concludere né \decision{RIMBORSABILE} né
\decision{NON\_RIMBORSABILE}. La lista dei campi mancanti accompagna sempre
la risposta, orientando il clinico verso la raccolta del dato necessario.

\subsection{Algoritmo di valutazione}

\Cref{alg:rule-eval} riporta il pseudocodice della funzione principale
\texttt{evaluate(caso, indice\_regole)}.

\begin{algorithm}[htbp]
\caption{Valutazione di un caso clinico (\texttt{evaluate})}
\label{alg:rule-eval}
\begin{algorithmic}[1]
\Require $caso$: dizionario campi clinici; $indice$: RuleIndex per la nota
\Ensure EvaluationResult con decisione, audit, flag clinici
\State $\mathit{vars} \gets \textsc{ComputaDerivedVariables}(caso)$
\State $\mathit{bypass\_set} \gets \emptyset$
\State \Comment{Fase EXCEPTION: identifica bypass prima di EXCL e INCLUSION}
\ForAll{$r \in indice.\text{exception\_rules}$ in ordine}
  \State $tv \gets \textsc{EvalCondizione}(r.\text{condition}, \mathit{vars})$
  \State $\textsc{ScriviAudit}(r, tv)$
  \If{$tv = \mathrm{TRUE}$}
    \State $\mathit{bypass\_set} \gets \mathit{bypass\_set} \cup r.\text{bypasses}$
  \EndIf
\EndFor
\State \Comment{Fasi SCOPE, EXCL\_HARD, INCLUSION, PATHWAY: pattern comune}
\ForAll{$\mathit{fase} \in [\textsc{Scope},\textsc{ExclHard},\textsc{Inclusion},\textsc{Pathway}]$}
  \State $\mathit{pending\_unknown} \gets \emptyset$
  \ForAll{$r \in indice.\mathit{fase}.\text{rules}$ in ordine}
    \If{$r.\text{id} \in \mathit{bypass\_set}$} \State \textbf{continua}
    \EndIf
    \State $tv \gets \textsc{EvalCondizione}(r.\text{condition}, \mathit{vars})$
    \State $\textsc{ScriviAudit}(r, tv)$
    \If{$tv = \mathit{fase}.\text{fail\_tv}$}
      \State \Return $\textsc{Finalizza}(\decision{NON\_RIMBORSABILE}, \mathit{audit})$
    \EndIf
    \If{$tv = \mathrm{UNKNOWN}$}
      \State $\mathit{pending\_unknown} \gets \mathit{pending\_unknown} \cup \{r\}$
    \EndIf
  \EndFor
  \If{$\mathit{pending\_unknown} \neq \emptyset$}
    \State \Return $\textsc{Finalizza}(\decision{NON\_DETERMINABILE}, \mathit{audit})$
  \EndIf
\EndFor
\State \Comment{Fasi di guidance: producono flag, non modificano la decisione}
\ForAll{$r \in indice.\text{guidance\_rules}$}
  \State $tv \gets \textsc{EvalCondizione}(r.\text{condition}, \mathit{vars})$
  \If{$tv = \mathrm{TRUE}$}
    \State $\textsc{AggiungiFlagClinico}(r)$
  \EndIf
\EndFor
\State \Return $\textsc{Finalizza}(\decision{RIMBORSABILE}, \mathit{audit})$
\end{algorithmic}
\end{algorithm}

Il routing fra le quattro Note avviene prima dell'ingresso nel ciclo: il campo
\texttt{note\_id} della richiesta seleziona l'oggetto \texttt{RuleIndex}
corrispondente, evitando qualsiasi ambiguità fra i cataloghi. Se il campo è
assente o non riconosciuto, l'orchestratore restituisce un errore 422 prima
di invocare il motore.

% ============================================================
\section{Pipeline RAG: ingestione e recupero}
\label{sec:rag-impl}

\subsection{Ingestione del corpus normativo}

Il modulo \texttt{rag\_pipeline/ingest\_v2.py} indicizza i PDF delle quattro
Note AIFA in collezioni ChromaDB dedicate. La versione~2 dell'ingestione
introduce il tracciamento degli offset di carattere all'interno di ogni
pagina, le bounding box per chunk, il checksum SHA-256 del testo e il
riconoscimento delle tabelle tramite \texttt{pdfplumber}.

L'estrazione del testo è eseguita con PyMuPDF per i blocchi testuali e con
\texttt{pdfplumber} per le tabelle, che vengono marcate con il flag
\texttt{is\_in\_table=True} e suddivise in un chunk per riga. Il testo puro è
tokenizzato in frasi con spaCy (modello italiano \texttt{it\_core\_news\_sm});
i chunk sono costruiti con una finestra scorrevole di dimensione target
\texttt{TARGET\_CHARS = 1800} caratteri e sovrapposizione di
\texttt{OVERLAP\_SENTS = 2} frasi, bilanciando la densità informativa con la
granularità necessaria alle ancore citazionali~\cite{gao2023ragsurvey}. I
chunk sono inviati a ChromaDB in lotti da 100 elementi con i metadati di
pagina, nota, offset e hash, usando il modello di embedding \embedmodel.

Una nota sull'esclusione deliberata: il file \texttt{nota-97-all-1.pdf} è
escluso dall'indicizzazione a causa di un problema di encoding Wingdings che
produce mojibake non recuperabile con le librerie disponibili. Il contenuto
normativo di Nota~97 è comunque coperto dai file \texttt{nota-97-all-2.pdf}
e \texttt{nota-97-all-3.pdf}, privi del difetto.

\Cref{alg:ingest} riporta il pseudocodice della funzione di indicizzazione
per una singola nota.

\begin{algorithm}[htbp]
\caption{Indicizzazione di una nota AIFA (\texttt{indicizza\_nota})}
\label{alg:ingest}
\begin{algorithmic}[1]
\Require $\mathit{pdf\_path}$: percorso del PDF; $\mathit{nota\_id}$: identificatore nota
\Ensure Numero di chunk indicizzati in ChromaDB
\State $\mathit{hash\_pdf} \gets \textsc{SHA256}(\mathit{pdf\_path})$
\State $\mathit{pagine} \gets \textsc{EstraiPagine}(\mathit{pdf\_path})$
\Comment{PyMuPDF rawdict + pdfplumber}
\State $\mathit{chunk\_list} \gets [\,]$
\ForAll{$p \in \mathit{pagine}$}
  \State $\mathit{blocchi} \gets \textsc{EstraiBlockiTesto}(p)$
  \State $\mathit{tabelle} \gets \textsc{EstraiTabelle}(p)$ \Comment{pdfplumber}
  \ForAll{$b \in \mathit{blocchi}$}
    \If{$|b.\text{text}| < \mathrm{MIN\_BLOCK\_CHARS}$} \State \textbf{salta}
    \EndIf
    \State $\mathit{frasi} \gets \textsc{TokenizzaFrasi}(b.\text{text})$ \Comment{spaCy it}
    \State $\mathit{finestra} \gets [\,]$; $\mathit{char\_count} \gets 0$
    \ForAll{$s \in \mathit{frasi}$}
      \State $\mathit{finestra}.\text{append}(s)$; $\mathit{char\_count} \mathrel{+}= |s|$
      \If{$\mathit{char\_count} \geq \mathrm{TARGET\_CHARS}$}
        \State $\mathit{chunk\_list}.\text{append}(\textsc{BuildChunk}(\mathit{finestra}, p, \mathit{nota\_id}))$
        \State $\mathit{finestra} \gets \mathit{finestra}[-\mathrm{OVERLAP\_SENTS}:]$
        \State $\mathit{char\_count} \gets \sum_{s \in \mathit{finestra}} |s|$
      \EndIf
    \EndFor
    \If{$\mathit{finestra} \neq [\,]$}
      \State $\mathit{chunk\_list}.\text{append}(\textsc{BuildChunk}(\mathit{finestra}, p, \mathit{nota\_id}))$
    \EndIf
  \EndFor
  \State \Comment{Tabelle: un chunk per riga}
  \ForAll{$\mathit{riga} \in \textsc{AppiattisciTabelle}(\mathit{tabelle})$}
    \State $\mathit{chunk\_list}.\text{append}(\textsc{BuildChunkTabella}(\mathit{riga}, p, \mathit{nota\_id}))$
  \EndFor
\EndFor
\State \Comment{Upsert in batch da CHROMA\_BATCH=100}
\State $\mathit{coll} \gets \textsc{GetOrCreateCollection}(\texttt{"nota\_"} + \mathit{nota\_id} + \texttt{"\_v2"})$
\ForAll{$\mathit{batch} \in \textsc{Partition}(\mathit{chunk\_list}, 100)$}
  \State $\mathit{coll}.\textsc{Upsert}(\mathit{batch})$
\EndFor
\State \Return $|\mathit{chunk\_list}|$
\end{algorithmic}
\end{algorithm}

\subsection{Recupero e reranking}

La fase di recupero invoca ChromaDB con una query di embedding della domanda
clinica e recupera i $k$ chunk più vicini per distanza coseno. I chunk
restituiti vengono poi riordinati da un cross-encoder, \reranker, che calcola
la rilevanza contestuale fra la query e ciascun testo in modo più accurato
rispetto alla similarità vettoriale~\cite{nogueira2019passage}. Solo i chunk
che superano la soglia di rilevanza del reranker vengono inclusi nel contesto
del prompt LLM, limitando il rischio di allucinazioni indotte da passaggi
normativi irrilevanti~\cite{ji2023hallucination}.

% ============================================================
\section{Orchestratore, riproducibilità e demo}
\label{sec:orchestrator-impl}

L'orchestratore è un servizio FastAPI esposto sulla porta 8001 che coordina
sequenzialmente il rule engine, il retriever semantico e il modello
linguistico. Il suo ciclo di vita è gestito dal context manager asincrono
\texttt{lifespan}, che al momento dell'avvio inizializza il
\texttt{CDSSOrchestrator} tentando prima il caricamento diretto del rule
engine in-process e ricorrendo alla comunicazione HTTP solo in caso di
fallimento, per eliminare la latenza di rete nei deployment integrati.

\subsection{Endpoint pubblici}

Il servizio espone due endpoint. L'endpoint \texttt{GET /health} restituisce
lo stato di inizializzazione dell'orchestratore e viene usato dai healthcheck
di Docker Compose per coordinare l'ordine di avvio dei container. L'endpoint
principale è \texttt{POST /explain}, autenticato tramite l'header
\texttt{X-API-Key}: la chiave attesa è letta dalla variabile d'ambiente
\texttt{AIFA\_API\_KEY} e la sua assenza rende il servizio aperto, facilitando
lo sviluppo locale senza configurazione aggiuntiva.

La struttura di input del \texttt{POST /explain} comprende il campo
\texttt{schema\_version} (valore atteso \texttt{3.3}), l'identificatore della
nota (\texttt{note\_id}: uno fra \texttt{"01"}, \texttt{"13"}, \texttt{"66"},
\texttt{"97"}), il farmaco (\texttt{drug\_id}) e due dizionari distinti:
\texttt{patient\_data} per i dati clinici misurabili e
\texttt{clinician\_asserted} per le asserzioni che richiedono giudizio clinico
diretto e non sono automaticamente derivabili da valori numerici. La
separazione dei due dizionari è una scelta progettuale deliberata: rende
esplicito al chiamante quali informazioni provengono da sistemi automatici e
quali richiedono l'intervento del medico.

La risposta \texttt{CDSSResponse} include la decisione di rimborsabilità
(\texttt{reimbursement\_decision}), la lista di flag clinici
(\texttt{clinical\_flags}), i chunk normativi recuperati
(\texttt{retrieved\_chunks}), la spiegazione testuale generata dall'LLM
(\texttt{explanation}) e le citazioni delle ancore PDF
(\texttt{citations}). Il motore non delega mai la decisione al modello
linguistico: l'LLM produce esclusivamente la spiegazione in linguaggio naturale
a partire da una decisione già fissata dal rule engine, in accordo con
l'architettura neuro-simbolica descritta nel \Cref{ch:progettazione}.

La gestione degli errori è progettata per non esporre i dati del paziente
nei messaggi di risposta HTTP: eventuali eccezioni sono tracciate server-side
nel log, mentre il client riceve soltanto un messaggio generico con codice
500, prevenendo la fuga accidentale di informazioni cliniche sensibili.

\subsection{Riproducibilità e cleanroom}
\label{subsec:reproducibility-impl}

La riproducibilità è garantita da tre meccanismi ortogonali. Il primo è il
file \texttt{requirements.lock} con hash esatti di tutte le dipendenze Python,
che impedisce aggiornamenti impliciti alle librerie. Il secondo è il manifest
\texttt{ingestion\_manifest\_v2.json}, che registra il checksum SHA-256 di
ogni PDF normativo ed è confrontato ad ogni avvio del servizio con i file
presenti su disco. Il terzo è la fissazione dei parametri di generazione del
modello linguistico a \texttt{temperature=0} e \texttt{seed=42}, che rende
deterministica la risposta testuale a parità di input e contesto. La pipeline
\texttt{run\_full\_overnight.sh} esegue la valutazione completa con
checkpoint per stage (ingestione, rule engine, risposte RAG, RAGAS,
QAEval) e può essere ripresa dal punto di interruzione in caso di riavvio.
La riproducibilità profonda è garantita dallo script
\texttt{full\_cleanroom\_v2.sh}, che cancella le collezioni ChromaDB,
reingerisce i PDF da zero, esegue la suite di test e valuta: solo questa
procedura, la \emph{cleanroom run}, esclude che i risultati siano
condizionati da artefatti di esecuzioni precedenti residui nel database
vettoriale, perché la valutazione sulla sola ChromaDB esistente può mascherare
bug della fase di ingestione che non emergono fino al successivo wipe.

\subsection{Suite di test e demo interattiva}
\label{subsec:test-demo-impl}

La suite conta \testcountall{} test verdi, organizzati in tre livelli. I test
unitari verificano il comportamento delle singole funzioni della logica
trivalente, degli operatori di confronto e del parser YAML. I test di
integrazione eseguono scenari clinici completi sul rule engine
(\testcountengine{} casi, di cui 131 specifici per le quattro Note AIFA) e
sull'orchestratore RAG (\testcountorch{} casi). I test \emph{property-based}
scritti con Hypothesis verificano invarianti come la
simmetria delle tavole di Kleene, la monotonia dell'esito al crescere
dell'informazione disponibile e la preservazione delle ancore normative in
tutte le risposte non vuote. La copertura complessiva del codice è
dell'85\%.

La demo interattiva, infine, è realizzata in Streamlit sulla porta 8501.
Streamlit è stata scelta per la possibilità di produrre un'interfaccia
clinicamente leggibile in poche centinaia di righe di Python puro: la minore
flessibilità di layout rispetto a un'applicazione React è accettabile per un
prototipo destinato alla revisione clinica e non a un deployment prescrittivo.
La sidebar espone in posizione fissa il disclaimer medico-legale obbligatorio;
il client incapsula le chiamate al rule engine e all'orchestratore in un
singolo modulo con gestione esplicita degli errori di connessione e
degradazione graceful in caso di indisponibilità di uno dei servizi. La
pagina di dettaglio include un pannello \emph{what-if} che consente di
modificare i parametri clinici e rieseguire la pipeline completa (rule engine
e generazione della spiegazione) aggiornando interattivamente sia la decisione
sia la motivazione; le identità dei pazienti fittizi sono generate
deterministicamente da seed fisso e gli avatar SVG sono prodotti localmente
senza richiamare servizi esterni, nel rispetto della riservatezza dei dati
anche nella fase dimostrativa.

Le scelte descritte in questo capitolo non sono fini a sé stesse: sono ciò che
rende le domande di ricerca verificabili. Il determinismo del rule engine e la
fissazione dei parametri di generazione sono il presupposto di RQ4, che misura
la stabilità dell'esito a ripetizioni e perturbazioni; la procedura di cleanroom
e il \texttt{requirements.lock} con hash garantiscono che i numeri con cui il
\Cref{ch:valutazione} risponde a tutte le RQ (\emph{Research Question}) provengano da una catena
riproducibile da PDF a risultato, e non da artefatti accumulati in esecuzioni
precedenti. È su questa base che il capitolo seguente affronta le sette domande
in modo quantitativo.
